% !TeX root = ../DECO_buku_iot_cerdas.tex
\chapter{PENDAHULUAN}

\section{Latar Belakang}

Perkembangan teknologi \textit{Internet of Things} (IoT) telah memungkinkan
berbagai sistem pemantauan berbasis sensor dikembangkan secara lebih mudah,
fleksibel, dan terhubung dengan platform \textit{cloud}. IoT memungkinkan
perangkat fisik seperti sensor, mikrokontroler, dan platform dashboard saling
berkomunikasi melalui jaringan internet untuk mengirim, menyimpan, dan
menampilkan data secara kontinu. Salah satu bentuk implementasi sederhana
namun relevan dalam IoT adalah sistem monitoring suhu berbasis sensor digital
\cite{thingspeak2024}, \cite{thingspeak_iot2021}.

Suhu merupakan salah satu parameter lingkungan yang penting untuk dipantau
pada berbagai aplikasi, seperti pemantauan ruangan, sistem penyimpanan,
laboratorium, perangkat elektronik, maupun lingkungan kerja. Sensor DS18B20
dapat digunakan untuk pengukuran suhu digital karena mendukung komunikasi
berbasis \textit{1-Wire} dan menyediakan pembacaan suhu dalam format digital
dengan resolusi yang dapat dikonfigurasi \cite{ds18b20datasheet}. Oleh karena
itu, sensor ini sesuai untuk digunakan pada sistem monitoring suhu berbasis
mikrokontroler.

Pada tugas besar ini, dirancang sistem monitoring suhu berbasis IoT
menggunakan ESP32-S3 dan sensor DS18B20. ESP32-S3 dipilih karena mendukung
konektivitas Wi-Fi 2,4 GHz sehingga dapat digunakan untuk mengirim data
sensor ke platform IoT berbasis internet \cite{esp32s3espressif}. Sistem
disimulasikan menggunakan Wokwi, yang menyediakan lingkungan simulasi
mikrokontroler serta koneksi internet virtual melalui gateway khusus untuk
proyek ESP32 \cite{wokwi_wifi}. Data suhu dari sensor DS18B20 dikirimkan oleh
ESP32-S3 ke platform ThingSpeak melalui koneksi Wi-Fi virtual Wokwi.

ThingSpeak digunakan sebagai platform IoT karena mendukung pengumpulan,
visualisasi, dan analisis data dari perangkat IoT secara
\textit{cloud-based} \cite{thingspeak2024}. Pada sistem ini, ThingSpeak
digunakan untuk menampilkan suhu aktual, prediksi suhu, dan status hasil AI
melalui dashboard berbasis grafik. Dashboard tersebut memenuhi kebutuhan
sistem IoT cerdas karena mampu menampilkan data secara otomatis dan
berkelanjutan sesuai dengan ruang lingkup tugas besar yang menekankan
dashboard monitoring, integrasi AI, dokumentasi, dan demonstrasi sistem
\cite{tugasbesar2026}.

Selain monitoring, sistem ini juga dilengkapi dengan kecerdasan buatan
berbasis \textit{Random Forest Regressor} yang dijalankan menggunakan Python
di Visual Studio Code. Random Forest merupakan metode \textit{ensemble
learning} yang menggabungkan banyak pohon keputusan untuk meningkatkan
kemampuan prediksi dan mengurangi risiko \textit{overfitting}
\cite{breiman2001}. Pada implementasi ini, model Random Forest membaca data
historis suhu dari ThingSpeak, kemudian memprediksi suhu 1 menit ke depan.
Model tersebut diimplementasikan menggunakan pustaka \textit{scikit-learn},
yang menyediakan modul \textit{RandomForestRegressor} untuk permasalahan
regresi \cite{sklearn_rf}, \cite{pedregosa2011}.

Dengan demikian, sistem yang dikembangkan tidak hanya berfungsi sebagai
sistem monitoring suhu, tetapi juga sebagai sistem IoT cerdas yang mampu
melakukan analisis prediktif sederhana. Implementasi ini mencakup akuisisi
data sensor, komunikasi IoT, dashboard monitoring, integrasi AI, dan
visualisasi hasil prediksi secara \textit{end-to-end}. Struktur laporan
disusun mengikuti format tugas besar yang memuat bagian pendahuluan,
arsitektur sistem IoT, kecerdasan buatan, dashboard dan komunikasi data,
serta simpulan \cite{laporan_contoh}.

\section{Rumusan Masalah}

Berdasarkan latar belakang tersebut, rumusan masalah pada tugas besar ini
adalah sebagai berikut:

\begin{enumerate}
    \item Bagaimana merancang sistem monitoring suhu berbasis ESP32-S3 dan
    sensor DS18B20 pada simulator Wokwi?
    \item Bagaimana mengirimkan data suhu dari Wokwi ke ThingSpeak secara
    periodik melalui koneksi Wi-Fi virtual?
    \item Bagaimana mengintegrasikan model AI \textit{Random Forest Regressor}
    untuk memprediksi suhu 1 menit ke depan berdasarkan data historis?
    \item Bagaimana menampilkan data suhu aktual, prediksi suhu, dan status
    AI pada dashboard ThingSpeak?
\end{enumerate}

\section{Tujuan Penelitian}

Tujuan penelitian pada tugas besar ini adalah:

\begin{enumerate}
    \item Merancang sistem monitoring suhu menggunakan ESP32-S3 dan sensor
    DS18B20 berbasis simulasi Wokwi.
    \item Mengintegrasikan ESP32-S3 dengan platform ThingSpeak untuk
    pengiriman dan visualisasi data suhu.
    \item Mengembangkan program AI berbasis \textit{Random Forest Regressor}
    menggunakan Python untuk memprediksi suhu 1 menit ke depan.
    \item Menampilkan hasil monitoring dan prediksi AI pada dashboard
    ThingSpeak.
    \item Mendokumentasikan sistem dalam bentuk laporan, \textit{source code},
    \textit{repository}, dan video demonstrasi sesuai ketentuan tugas besar
    \cite{tugasbesar2026}.
\end{enumerate}

\section{Potensi dan Segmentasi}

Penelitian ini memiliki potensi teknis dan edukatif dalam pengembangan sistem
IoT cerdas berbasis pemantauan suhu. Penggunaan sensor DS18B20 bersama
ESP32 memberikan pendekatan akuisisi data yang hemat daya dan dapat
diterapkan pada berbagai perangkat IoT berdaya rendah. Integrasi model
\textit{Random Forest Regressor} pada sistem memperlihatkan bahwa analisis
prediktif dapat dilakukan secara eksternal pada komputer tanpa membebani
sumber daya mikrokontroler \cite{breiman2001}.

Dari sisi platform, ThingSpeak menyediakan API yang mudah diintegrasikan
dengan perangkat IoT maupun skrip Python, sehingga visualisasi data dan
pengiriman hasil AI dapat berjalan dalam satu ekosistem \textit{cloud}
yang sama \cite{thingspeak2024}. Hasil penelitian ini diharapkan menjadi
referensi implementatif bagi pengembangan sistem monitoring berbasis IoT
cerdas di lingkungan akademik.

Segmentasi penelitian ditujukan untuk lingkungan akademik, mahasiswa yang
mempelajari IoT dan \textit{machine learning}, serta pengembang sistem
pemantauan kondisi lingkungan berbasis \textit{cloud}. Fokus penelitian
berada pada demonstrasi integrasi end-to-end antara perangkat IoT,
platform cloud, dan model AI prediktif, bukan pada optimasi akurasi
model secara mendalam.

\section{Batasan Masalah dan Batasan Penelitian}

Agar sistem tetap fokus dan realistis, batasan pada tugas besar ini adalah:

\begin{enumerate}
    \item Sistem menggunakan satu sensor suhu digital DS18B20.
    \item Mikrokontroler yang digunakan adalah ESP32-S3 pada simulator Wokwi.
    \item Platform IoT yang digunakan adalah ThingSpeak.
    \item AI dijalankan di komputer menggunakan Python pada Visual Studio
    Code, bukan langsung pada ESP32-S3.
    \item Model AI yang digunakan adalah \textit{Random Forest Regressor}.
    \item Prediksi suhu dilakukan untuk \textit{horizon} waktu 1 menit ke
    depan (4 langkah prediksi dengan interval 15 detik per data).
    \item Dashboard ThingSpeak menggunakan tiga \textit{field} utama, yaitu
    suhu aktual, prediksi suhu, dan status AI.
    \item Sistem tidak menggunakan aktuator seperti LED, \textit{buzzer},
    \textit{relay}, atau motor.
\end{enumerate}

\section{Metode Penelitian}

\subsection{Studi Literatur}
Tahap ini dilakukan dengan mempelajari referensi mengenai IoT, sensor DS18B20,
mikrokontroler ESP32, simulator Wokwi, platform ThingSpeak, metode
\textit{Random Forest Regressor}, serta metrik evaluasi model regresi
\cite{breiman2001}, \cite{sklearn_rf}, \cite{pedregosa2011}.

\subsection{Analisis Kebutuhan dan Perancangan Sistem}
Tahap ini meliputi perancangan arsitektur sistem IoT secara \textit{end-to-end},
penentuan komponen perangkat keras dan perangkat lunak, perancangan wiring
sensor DS18B20 pada ESP32, konfigurasi channel ThingSpeak, serta perancangan
fitur dan alur kerja model AI berbasis \textit{Random Forest Regressor}.

\subsection{Implementasi Sistem}
Sistem diimplementasikan menggunakan simulator Wokwi untuk bagian perangkat
keras (ESP32 dan DS18B20) dan Python pada Visual Studio Code untuk bagian AI.
Implementasi mencakup program pembacaan sensor, pengiriman data ke ThingSpeak,
pembentukan fitur \textit{time series}, pelatihan model, prediksi suhu, dan
pengiriman hasil prediksi kembali ke ThingSpeak.

\subsection{Pengujian dan Evaluasi Sistem}
Pengujian dilakukan dengan menjalankan simulasi Wokwi dan program Python secara
bersamaan. Data suhu aktual, prediksi suhu, dan status AI diamati melalui
dashboard ThingSpeak. Evaluasi model dilakukan menggunakan metrik MAE, RMSE,
dan $R^2$ \textit{Score} pada data uji.

\subsection{Analisis dan Penarikan Kesimpulan}
Data hasil pengujian dianalisis untuk menilai keberhasilan setiap komponen
sistem, akurasi prediksi model AI, dan konsistensi klasifikasi status suhu.
Hasil analisis digunakan untuk menjawab rumusan masalah dan menyusun simpulan
serta saran pengembangan sistem ke depan.
